\section{Seguint la pista de l'energia}

Hem parlat que l'energia no desapareix: es transforma i passa d'un lloc a un altre. Ara intentarem seguir-li la pista en un experiment real.

Bona part de l'energia que utilitzem a la Terra prové, directament o indirectament, del Sol. En un aerogenerador, per exemple, l'energia del moviment de l'aire es transforma en energia elèctrica. Però aquell vent tampoc ha aparegut del no-res: l'escalfament desigual de l'atmosfera per part del Sol, juntament amb la rotació de la Terra, hi té molt a veure.

L'energia té aquesta mania: costa força trobar-la quieta.

En el nostre experiment, l'energia arriba en forma d'electricitat a una bombeta. El corrent passa per un filament que ofereix molta resistència i aquest s'escalfa. No és exactament el mateix mecanisme que quan et fregues les mans, però el resultat visible s'hi assembla: part de l'energia acaba convertida en energia tèrmica.

Quan el filament s'escalfa prou, comença a emetre \textbf{radiació electromagnètica}. El nom pot sonar una mica més seriós del compte: físicament estem parlant de la mateixa família de fenòmens a la qual pertany la llum que veiem.

Una part d'aquesta energia arriba al sensor de la càmera i es transforma en senyals elèctrics i, finalment, en informació digital. Una altra part escalfa l'aire, la bombeta, el suport i fins i tot una mica la caixa.

Seguir la pista de tota l'energia pot arribar a ser una feina bastant llarga.

\subsection{Activitat 1. On ha anat a parar l'energia?}

Fes un esquema del nostre experiment i intenta representar el camí que segueix l'energia, però no et quedis només amb el tros que veiem damunt la taula.

Pots començar per:
\[
\text{energia elèctrica}
\longrightarrow
\text{calor}
\longrightarrow
\text{radiació}
\longrightarrow
\text{sensor}
\]

Ara mira una mica més lluny: què ha hagut de passar abans perquè aquesta energia arribés a la font d'alimentació? I què passa després que la llum arribi a la càmera?

Completa l'esquema indicant altres transformacions o llocs on creguis que també acaba part de l'energia. Nosaltres només observem una part de la història.


\subsection{Activitat 2. Comencem a donar-li volts}

Ara observarem la bombeta amb la imatge en color de la càmera.

Comença amb una tensió baixa i augmenta-la progressivament. Anota el \textbf{voltatge} a partir del qual la llum comença a ser visible. Per cert, \textbf{tensió} i \textbf{voltatge} són dues maneres de parlar de la mateixa cosa.

Observa quan comença a aparèixer llum i com va canviant el seu color a mesura que augmentes la tensió.

\subsection{Activitat 3. Separant els colors}

La càmera construeix la imatge en color a partir de tres canals:
\[
R \text{ (red)} \qquad G \text{ (green)} \qquad B \text{ (blue)}
\]

és a dir:
\[
\text{vermell} \qquad
\text{verd} \qquad
\text{blau}.
\]

Ara els observarem per separat.

Augmenta lentament la tensió i apunta quan comences a detectar clarament cadascun dels tres canals.
\begin{center}
\renewcommand{\arraystretch}{1.6}
\begin{tabularx}{\linewidth}{|X|
                                >{\centering\arraybackslash}X|
                                >{\centering\arraybackslash}X|}
\hline
\textbf{Canal} &
\textbf{Tensió (V)} &
\textbf{Intensitat (A)} \\
\hline
Vermell & & \\ \hline
Verd    & & \\ \hline
Blau    & & \\ \hline
\end{tabularx}
\end{center}

Ara compara els tres canals aproximadament a \textbf{5 V} i després a \textbf{9 V}.

Quin canal sembla dominar quan el filament encara està relativament fred?

I quan augmentem la tensió?

Els tres colors van apareixent de manera completament aleatòria o sembla que hi ha un cert ordre?

Si creus que hi ha un ordre, intenta escriure'l:
\[
\boxed{\hspace{1.5cm}}
\longrightarrow
\boxed{\hspace{1.5cm}}
\longrightarrow
\boxed{\hspace{1.5cm}}
\]

Potser has vist alguna vegada una peça de ferro molt calenta. Primer es torna vermellosa i, si continuem escalfant-la, el color va canviant.

No sembla una coincidència.

\subsection{Activitat 4. I si encara hi hagués més?}

Fins ara hem separat la llum que podem veure en vermell, verd i blau.

Però aquí apareix una pregunta bastant natural:
\begin{quote}
Si aquests colors semblen seguir un cert ordre quan augmenta la temperatura, podria continuar aquest ordre més enllà dels colors que els nostres ulls poden veure?
\end{quote}

Sabem que hi ha animals capaços de detectar radiació que nosaltres no veiem. També hi ha càmeres capaces de fer-ho.

La nostra n'és una.

Abans d'activar-la, fes una predicció:
\begin{quote}
Si existeix una radiació que apareix abans que el vermell visible, esperaries detectar-la amb una tensió més baixa o més alta?
\end{quote}

Ara activa el canal \textbf{infraroig} de la càmera.

L'infraroig forma part de la mateixa radiació electromagnètica que la llum visible, però queda fora del rang que poden detectar els nostres ulls.

Augmenta de nou la tensió des d'un valor baix i compara quan apareix amb el canal vermell.
\begin{center}
\renewcommand{\arraystretch}{1.6}
\begin{tabularx}{\linewidth}{|X|
                                >{\centering\arraybackslash}X|
                                >{\centering\arraybackslash}X|}
\hline
\textbf{Canal} &
\textbf{Tensió (V)} &
\textbf{Intensitat (A)} \\
\hline
Infraroig & & \\ \hline
Vermell   & & \\ \hline
\end{tabularx}
\end{center}

Quin apareix abans?

Ara puja fins a la tensió a la qual començaves a veure clarament el vermell.

Compara els dos canals:
\[
\text{infraroig}
\qquad\text{vs.}\qquad
\text{vermell}
\]

Quin sembla més intens en aquell moment?

\subsection{Discussió. La llum és més gran que els nostres ulls}

Ara ja podem tornar al principi.

La llum visible és només una petita part de la radiació electromagnètica. Quan diem ``radiació'' no estem parlant necessàriament d'alguna cosa estranya o perillosa: la llum que veiem també és radiació electromagnètica.

Els objectes, pel simple fet de tenir temperatura, emeten radiació. A temperatures relativament baixes, bona part queda fora del que poden veure els nostres ulls. Quan augmentem molt la temperatura, l'emissió va entrant progressivament dins del visible.

Per això una peça de metall, una brasa o el filament d'una bombeta poden començar a brillar.

I l'ordre que hem observat continua més enllà del que poden veure els nostres ulls. Nosaltres només detectem una petita part de tota la radiació que existeix. Altres animals tenen una finestra una mica diferent. I les càmeres, si les construïm per fer-ho, també.

Pots comprovar aquesta idea buscant informació en un documental, un llibre o qualsevol altra font fiable.


\subsection{Una última idea}

L'energia es transforma constantment.

Sabem que l'energia és una magnitud física de la qual podem mesurar els efectes i les transformacions. En aquest experiment l'hem seguida mentre passava per electricitat, escalfament, radiació i informació.

I segurament encara ens hem deixat coses pel camí.
